\section{Introduction}
For a finite set $U\subset\R^2$, let
\[
 u(U)=\#\{\{x,y\}\subset U:|x-y|=1\},\qquad
 u(n)=\max_{|U|=n}u(U),
\]
where the pairs are unordered and the distance is Euclidean.
Erd\H{o}s introduced the unit distance problem in 1946
\cite{Erdos1946}. His lattice construction gives
\[
 u(n)\ge n^{1+c/\log\log n}
\]
for some $c>0$ and arbitrarily large $n$. The upper bound
\[
 u(n)=O(n^{4/3})
\]
is due to Spencer, Szemer\'edi and Trotter
\cite{SpencerSzemerediTrotter1984}, with a later crossing-number proof
due to Sz\'ekely \cite{Szekely1997}.

The recent number-field construction of OpenAI \cite{OpenAI2026},
explained by Alon et al.~\cite{Alon2026}, gives
$u(n)\ge n^{1+\delta}$ for a fixed $\delta>0$ and arbitrarily large $n$.
Sawin \cite{Sawin2026} made this exponent explicit. In our convention
for unordered pairs, his result gives
\[
 u(n)\ge C^{-1}n^{1.014114}
\]
for an absolute constant $C$ and arbitrarily large $n$.
Emmerich \cite{Emmerich2026} obtained a certificate with
$\delta=0.015263\ldots$ by optimizing the integer parameters in Sawin's
criterion. A subsequent certificate of Emmerich and Cordella
\cite{EmmerichCordella2026} reports $\delta=0.031185334443\ldots$.
The author's earlier manuscript \cite{NaslundLegacy2026} records
$\delta=0.0358324$. We compare its construction with the present one below.

Our main result is:
\begin{theorem}\label{thm:main}
There are finite sets $U_j\subset\R^2$, with $|U_j|\to\infty$, such that
\[
 \frac{u(U_j)}{|U_j|^{1.0418235}}\longrightarrow\infty.
\]
In particular, $u(n)\ge n^{1.0418235}$ for arbitrarily large integers $n$.
\end{theorem}

We obtain the improvement by changing the number fields and counting
norm-one unit orbits in mixed signature, then optimizing the local and
archimedean overlap in the same construction. The result concerns an
unbounded sequence of cardinalities. No generalized Riemann hypothesis
or asymptotic class-number equality is used.

\subsection{Sawin's Construction}

The arithmetic source of unit distances is a quadratic extension
$K/F$ with involution $\iota$. If a prime $\mathfrak p$ of $F$ splits as
$\mathfrak P\iota\mathfrak P$, then the $k+1$ ideals
\[
 \mathfrak P^j(\iota\mathfrak P)^{k-j},\qquad 0\le j\le k,
\]
have the same relative norm $\mathfrak p^k$. Choices at different primes
multiply. An ideal-class argument converts some of these choices into
elements of a common relative norm. At an embedding where $\iota$
acts as complex conjugation, their images have a common length.
After rescaling, they give unit displacements in a planar projection
of an ideal lattice.

Sawin already uses fractional ideals, relative class numbers, and
powers of Frobenius elements to obtain small residue degrees
\cite{Sawin2026}. His explicit construction takes $F$ totally real and
$K/F$ CM. A sup-norm window controls the number of vertices and the
overlap under unit displacements, while Louboutin's relative
class-number bound \cite{Louboutin2000} controls the loss in selecting
an ideal class. Prescribing the local behavior requires relations in
a pro-$2$ group. The Golod--Shafarevich criterion
\cite{GolodShafarevich} ensures that sufficiently few such relations
leave an infinite tower. The use of Frobenius cuts to control splitting
belongs to the class-field-tower methods of Hajir, Maire and Ramakrishna
\cite{HajirMaireRamakrishna2021Cutting}.

There are two related arithmetic precedents. Small split primes and
ideal-class selection occur in the work of Ellenberg and Venkatesh
\cite{EllenbergVenkatesh2007} on class-group torsion. Number fields
with controlled discriminants have also been used to construct codes
and high-dimensional sphere packings
\cite{Lenstra1986,LitsynTsfasman1987}. Here the required lattice
vectors must have the same length after projection to the plane,
which is supplied by their common relative norm.

\subsection{The Earlier Bound $\delta=0.0358324$}

The earlier manuscript \cite{NaslundLegacy2026} keeps the totally
real/CM setting but changes several estimates in Sawin's criterion.
For Euclidean balls of radius $R>1/2$ and a unit vector $e_1$, the
overlap has exponential rate
\[
 \frac{1}{2d}\log
 \frac{\operatorname{vol}_{2d}(B_R\cap(B_R-e_1))}
      {\operatorname{vol}_{2d}(B_R)}
 \longrightarrow \frac12\log\left(1-\frac1{4R^2}\right).
\]
This replaces the loss $\log(1-1/R)$ in the earlier window estimate.
A bound for the dual lattice, together with Banaszczyk's transference
theorem \cite{Banaszczyk1993}, replaces the packing estimate for the
number of vertices by a covolume estimate. The narrow ideal-class
calculation also removes an exponential sign loss.

The analytic estimate uses a fixed multiquadratic subfield of every
field in the tower. Its normalized Euler product retains more
information than a comparison with a power of the Riemann zeta
function. This develops the small-prime comparisons in relative
class-number estimates \cite{Louboutin2000,Louboutin2005}.
The local construction allows dyadic ramification and charges some
relations at degree three in the filtration. For its set of the first
38 odd primes, the displayed infinitude certificate is
\[
 P(t)=1-38t+355t^2+114t^3,\qquad
 P(6/115)=-\frac{101}{1520875}<0.
\]
These changes, followed by reoptimization, give the earlier manuscript's
stated exponent $1.0358324$.

Some of these ideas were described in the author's MathOverflow
answer \cite{NaslundMO2026}, which credits spiderduckpig with placing
the second split-prime Frobenius relation in filtration degree at least
three. Sawin's comment there
also records that he knew the fixed-subfield analytic comparison,
without having evaluated it. The answer and the archived manuscript
are distinct records. In particular, the answer's June 9 change log
reports a bookkeeping correction and the figure $\delta>0.0357$.
We use the archived manuscript's $0.0358324$ to identify the earlier
construction, rather than as a hypothesis of the present proof.

The further improvement comes from allowing a different signature
and different local groups. Covolume counting, ball overlap and the
removal of the sign loss are already part of the earlier argument.
The new counting problem is to use a positive-rank group of norm-one
units without losing its regulator.

\subsection{Mixed Signature and Norm-One Units}

Let $K/\Q$ be totally imaginary and Galois, choose actual complex
conjugation $\iota$, and set $F=K^{\langle\iota\rangle}$. Write
\[
 [K:\Q]=2d,\qquad [F:\Q]=d=b+2c,\qquad \theta=c/d,
\]
where $(b,c)$ is the signature of $F$. The chosen complex embedding of
$K$ restricts to a real embedding of $F$ and supplies the planar
projection. The other embeddings
need not be real, since $\iota$ need not be central. In fact,
\[
 \frac bd=\frac1{|\operatorname{class}(\iota)|},
\]
where the conjugacy class is taken in $\Gal(K/\Q)$. Retaining a finite
quotient with many conjugates of this involution gives a field whose
signature is predominantly complex.

In the CM case, a norm-one element has modulus one at every
archimedean coordinate and the norm-one unit group is finite.
Above a complex place of $F$, its two coordinates can instead have
moduli $e^u$ and $e^{-u}$. The norm-one units have rank $c$ and
translate these reciprocal pairs. We average over a fundamental
domain of their full logarithmic lattice. Tiling gives an exact
integral divided by the relative regulator.

For the extensions used here, which are unramified at every finite
place, the class-number formula cancels this regulator together with
the kernel of extension of ideal classes. More precisely, if
$h_{\rm rel}=h_K/h_F$, $\kappa$ is the order of that kernel,
$R_U$ is the norm-one unit covolume, $w_K$ counts roots of unity, and
$\lambda=\rd(K)$, Proposition~\ref{geo:mass} gives
\[
 \frac{w_K}{h_{\rm rel}\kappa R_U}
 =\frac{2^{d-1}\pi^{b+c}}{\lambda^{d/2}L_F(1)},
 \qquad L_F(s)=\frac{\zeta_K(s)}{\zeta_F(s)}.
\]
At $s=1$, the quotient denotes the quotient of residues.
The averaging uses the entire unit lattice, so no independence of
its coordinates or equidistribution of unit orbits is needed.

\subsection{Local Relations and the Euler Calculation}

To construct the fields, we use the filtered images of the local
relation modules. Fox derivatives \cite{Fox1953} retain the
dependencies among relators that a count of their degrees alone
would discard. The filtration is the one in the
Golod--Shafarevich method, with its group-algebra description
\cite{Jennings,Quillen,Ershov2012}. In particular, a relation shared
by a global condition and a local condition is subtracted only after
identifying the same row in their actual filtered images.

This permits a nonabelian decomposition group of order $32$ at $2$,
with inertia order $8$ and residue degree $4$. Its relation calculation
allows residue degree $2$ at both $3$ and $5$. The resulting family has
\[
 \lambda=2^{9/4}\sqrt{15015},\qquad
 \theta\ge\frac{4095}{8192}.
\]
The gain at $3$ and $5$ has costs. Compared with the preceding
five-prime family, the root discriminant increases from
$4\sqrt{15015}$, and the residue degree at $17$ increases from $2$
to $4$. The smaller residue degrees also increase the corresponding
absolute Euler factors. The exponent improves because the additional
norm choices outweigh these losses in the final inequality.

For the analytic calculation we retain two fields
\[
 N\subset M\subset K,\qquad
 [N:\Q]=16384,\qquad [M:\Q]=524288.
\]
The full Euler calculation is made in $N$. Its regular representation
contains $128$ linear, $256$ quadratic, $456$ quartic and $124$ octic
irreducible representations, and we evaluate their Hecke factors
with the required multiplicities. Further Frobenius information from
$M$ is charged only on the remaining primes. This is necessary because
adding savings from different subfields can otherwise count the same
prime more than once.

The evaluation is at $\sigma=12001/12000$. The unconditional
Tsfasman--Vl\u{a}du\c{t} prime inequality \cite{TsfasmanVladut2002}
then gives, for all sufficiently large fields in the tower,
\[
 \frac1d\log L_F(1)<0.042161819.
\]
Thus the fixed-subfield idea from the earlier paper is used with a
nonabelian field and a complete factor calculation. The finite
field identities, conductor data and analytic remainders are proved
in sections \ref{ar:field} and \ref{an:section}.

\subsection{Archimedean Profiles and Finite Shells}

At a complex place of $F$, a norm-one displacement couples two
coordinates through reciprocal scaling. We consequently use a
profile on $\C^2$ that depends jointly on their radii. It is a positive
cubic Bernstein modification of a Student profile. The coordinates
above real places use a Gaussian. Products are taken over places,
while the two coordinates at each complex place remain coupled.

At a split finite place of residue-field size $Q$, a hard valuation
window gives the factor
\[
 (k+1)Q^{-\delta k}.
\]
We replace it by a weighted sum of six valuation shells. The shell
overlap matrix is explicit, so the gain and the mass can be optimized
together. Section \ref{fw:section} gives the exact finite functional.

Both choices enter the same arithmetic and geometric inequality.
For hard finite windows its form is
\[
 \sum_q\frac{\log(k_q+1)-\delta k_qf_q\log q}{e_qf_q}
 -(1/2-\delta)\log\lambda-C+\mathcal A_\delta(\theta)>0,
\]
where $e_q,f_q$ are absolute local indices, $C$ bounds
$d^{-1}\log L_F(1)$, and $\mathcal A_\delta$ is the archimedean
contribution. Proposition~\ref{geo:transfer} gives the precise
expression, and Proposition~\ref{fw:transfer} replaces the hard
windows by shell weights. Changing the local fields changes several
terms at once, which is why the tower and the profiles are optimized
together.

The weights are converted into an ordinary finite set by restricting
the combined energy to a concentration window. Ideal-class selection
uses the resulting overlap weights. A uniform Poisson estimate then
bounds the vertices for the same translate and unit deformation that
supplies the edges. At
\[
 \delta=\frac{83647}{2000000}=0.0418235,
\]
the rational witness in section \ref{sec:certificate} leaves a
certified margin greater than $5.33204359816794\times10^{-6}$
after the concentration allowance.

\subsection{Computation and Formalization}

The proof uses finite linear algebra, exact coefficient recurrences,
prime data and directed interval enclosures. The arithmetic inputs
and the algorithms for checking them accompany the manuscript.
Section \ref{sec:certificate} distinguishes computations rerun by
the supplied verifier from stored finite data. These certificates
concern numerical inequalities and explicit arithmetic, rather than
the statistical accuracy of a floating-point approximation.

There is also a separate Lean development, using Lean~4 and Mathlib
\cite{Lean4_2021,Mathlib2020}. Its selected theorem proves the planar
conclusion from one explicit inequality for the ordinary Dedekind
zeta function of the fixed field $M$. That inequality is stated in
Remark~\ref{an:fixed-field-form}. It has not been proved in Lean.
The formalization is consequently conditional, and is not used as a
substitute for the analytic argument or finite certificate here.
